\chapter*{Abstract}
\addcontentsline{toc}{chapter}{Abstract}
Photographs can reveal visible corrosion without establishing either internal damage or remaining service life. This technical activity examines that distinction in ferrocement using a predecessor dataset of 48 specimens from two experimental campaigns. The current aligned basis contains 791 readable photographic observations, but wire-area loss and ultimate load are each measured only once per specimen, at terminal testing. Campaign identity is aligned with mesh configuration, chloride concentration, and terminal exposure duration.

Saved experiments are synthesized by scientific question and independently checked through their tables, predictions, and split manifests. Classical colour features, frozen residual-network embeddings, and interpretable image descriptors support visible-corrosion estimation under specimen-grouped evaluation. Some strong results are label-adjacent: total-rust coverage and a rust-area feature have Spearman correlation 0.9996. This limits their interpretation as independent predictive discovery.

The focused terminal-load experiment compares image-only, metadata-only, and combined feature families. A metadata-only Ridge model has mean terminal-test MAE of 0.173\,kN, with fold standard deviation 0.046\,kN, across ten folds covering 48 specimens. Adding image descriptors does not establish a reliable pooled improvement. A post-onset sensitivity analysis has a limited gain in MAE for a combined configuration, without uniformly better ranking. Within each mesh family, every selected feature-set model has negative mean test $R^2$. Campaign exclusion is a severe extrapolation stress test, not an ordinary estimate of deployment performance.

Model-derived structural trajectories support only exploratory threshold-status analysis. The refined configuration has zero future threshold crossings, while crossings at baseline and during observation remain present. No lifetime-event supervision exists. The separate four-class severity branch has completed data preparation but no trained model. The principal contribution is therefore a bounded account of what the observations and evaluation design support, together with a concrete path toward stronger validation.

\chapter*{Executive Summary}
\addcontentsline{toc}{chapter}{Executive Summary}
\textbf{Scientific conclusion.} Surface assessment is the most defensible image-based task in this dataset. The available experiments do not establish transferable structural-condition inference or remaining-life prediction. This is a result about the present observations, representations, and evaluation procedures, not an impossibility claim about imaging.

\textbf{Why the distinction matters.} A photographed surface, a destructively measured wire cross-section, and a terminal bending load are different quantities. Repeated images provide rich visual supervision; repeating a specimen's terminal target alongside those images does not create repeated structural measurements. A model can exploit real design information and campaign-specific associations simultaneously.

\textbf{What supports the conclusion.} The most mature load analysis uses a complete feature-family comparison. Its pooled metadata-only result is competitive with all image-enhanced alternatives. The loss of performance within a mesh family and across campaigns limits the interpretation of the pooled score. Surface-treatment transfer is phase-dependent: one interpretable representation deteriorates under treatment exclusion, while a later representation does not show the same pattern.

\textbf{What remains useful.} The work provides reproducible aggregation of saved results, specimen-disjoint evaluation manifests, a clear separation between observed and derived quantities, and a prepared classification dataset. Negative results identify precisely which claims require new evidence.

\textbf{Next research decisions.} Complete the four-class classifier as a separate, controlled experiment; repair current configuration and source defects before retraining; evaluate early-image prediction at predeclared time horizons; and collect crossed experimental designs and independent structural or lifetime outcomes for stronger structural and prognostic claims. The report rebuild itself performs no model training.
